% !TeX root = ../papers/nc-p3.tex
\section{片側反転への全流入と二成分の交点の局所環}
\label{sec:mixed-flip-slice}
\begin{theorem}\label{thm:mixed-flip-slice}
半単純直線の片側反転 $s=\iota$, $L=T(a,1)A$ を閉軌道とする二切断階数スキームを考える。
$a^4\ne1$ では、全基底群の半安定商のこのファイバーは一軌道だけである。
$a^2=1$ では閉軌道と、定理\ref{thm:split-two-eigenvalue-jordan}の例外Jordan窓の
ちょうど二軌道で尽きる。双方とも最小周期1である。
この閉点での軌道横断面の完備局所環は
\[
                   k[[z,x,y]]/(zx,zy).
\]
直線が片側反転成分、平面が定理\ref{thm:split-simple-surface}の追加成分である。
二つは軌道方向を除いて一点だけで交わり、埋込み成分はない。
残る $a^2=-1$ の横断接表現は自明指標一つと、下記の相異なる七指標各一つである。
\end{theorem}
\begin{proof}[全パラメータの横断接指標]
片側反転窓で最初の色のビットを位置2,3だけで反転し、この色を二本の文字 $x_0,x_1$ に分ける。
他方の二文字を $Y$ とする。全非零語に両 $x$ が一つずつ現れる。
二つを独立に拡大するトーラスは射影安定群の恒等成分である。
実際、定理\ref{thm:all-mixed-flip-component-regular}の軌道微分の全値での階数から
恒等成分の次元は2である。語の横断指標は $(n_0-1,n_1-1)$、
$n_j$ は $x_j$ の出現数である。

二切断の接方程式をこの指標ごとに $\mathbb Q(\iota)[a,a^{-1}]$ で取る。
\path{computations/verify_mixed_flip_slice.py} は純関係に合わせて第二色のビットも奇位置で反転し、
各辺の六降順対をピボットにした左右の核から200積を作る。分母は $a$ の冪と定数だけである。
Smith形の一般核次元は和60、自明指標で20である。
非単元不変因子の全ての根と、各根で一次元ずつ増える指標は
\[
\begin{array}{c|l}
a^2=1 & (1,-1),(-1,1)\\
a^2=-1 & (-1,0),(0,-1),(-1,2),(2,-1),(0,2),(1,1),(2,0)
\end{array}
\]
である。59次元のアフィン基底軌道のうち自明指標の次元は19である。
実際、各位置の同重みブロックのLie環次元は $1+1+4=6$、
四位置の線形逆像でアフィン安定群の次元は $64-59=5$ で、全て重み零に入る。
他の指標でも軌道次元は一定である。よって横断接の全指標は、一般の自明指標一つと
表の該当箇所の和である。

閉軌道でLuna横断面を取り、安定群の極大トーラスで点へ流入する部分を考える。
量子三表の証明と同じく、強エタールな飽和開近傍での商の基底変換により
この基点のファイバー全体が横断面の零ファイバーで得られる。
接余接の指標座標を同変に持ち上げれば、それらは局所環の完備化を位相的に生成する。
一径数で基点に流れる弧では重み零・負の接座標は全て零であり、使える座標は正のものだけである。
$a^4\ne1$ では何もなく、閉点だけを得る。
$a^2=1$ ではどの余指標にも正となる接座標は高々一つで、二つの符号は互いに反対である。
\end{proof}

\begin{proof}[横断的な実現と全局所式]
同型な片側反転表示に移り、$E=(1,X,Z,XZ)$、$D_0=(Z,1,Z,1)$ とする。
第一対応は恒等、第二対応は $ED_0$ とし、第一交代テンソルを
\[
 P=P(\iota)+x(E_{0011}+\iota E_{1001})+
                  y(E_{1100}+\iota E_{0110})
\]
とする。原点の近傍ではこの全平面で両Schur式が成り立つ。
実際、追加曲面表示 $(1,-\iota v,\iota,v)$ の最初の係数を1に直すと、
上の二括弧の係数が共に $(1+\iota v)/(1-\iota v)$ になる。
$T(1,d)$ がこの係数を $d^{-1},d$ 倍し四対応を保つため、$xy\ne0$ の稠密な点で表示を得る。
原点でも両隣接逆行列の分母は単元で、Schur式は平面上で延びる。
座標軸の非零点は第二対応の片側反転を保つ二鎖Jordanであり、前記の例外一類になる。
一方、$x=y=0$ のまま $L=T(1,a)ED_0$ を動かすと片側反転曲線が得られる。
$z=a-1$ をその局所座標とする。

三接方向が軌道横断で独立で、混合二次障害 $zx,zy$ も独立であることを示す。
上の検算は原点のテンソルを二つのSchur式で再構成し、平面と曲線を各座標で微分する。
法109で $\iota=33$ とし、64個の基底行列単位の列に順に $z,x,y$ の微分を加えた階数は
$59,60,61,62$ である。各切断行列 $M$ の単元ピボット上の一般逆を $K$、
左右核を $L_0,R_0$ とし、二つの一次変分を $S,T$ とする。
Schur式の混合二次係数を消すには、二切断を合わせた接方程式の像に
$L_0(SKT+TKS)R_0$ がなければならない。
全接行列の階数は上で標数零でも194であり、法109で $zx,zy$ の右辺を順に加えた階数は
$194,195,196$ である。よって標数零の余核で二列は独立である。

滑らかな軌道座標を除くには、周囲の射影空間において軌道接の補空間と群の局所座標で
形式同型を取り、上の曲線・平面を横断面に射影する。
接が独立なので、三つの極小形式座標の中でこの二つの滑らかな像を座標直線と座標平面に同時に直せる。
実際、曲線と平面の各パラメータ表示の和から共通原点を引いたものは可逆な一次項を持つ。
従って横断面のイデアルは $(zx,zy)$ に含まれる。他方、接行列の独立な行から
余分な周囲座標を形式的に消去した後、残る方程式の二次項は上記の余核での障害である。
独立な二式が $(zx,zy)/\mathfrak m(zx,zy)$ を張るので、中山の補題からイデアルは全体に等しい。
これで完備環と二成分への所属を得る。

二つの正指標方向の各々には実際に平面の座標軸上のJordan表示がある。
横断面の同じ正接座標を持つ流入点は一意である。実際、トーラスで規格化して値を合わせ、
流入弧に沿って他の全座標を完備環の冪級数として比較すれば二弧は一致する。
非零値どうしはトーラスで移り、二軸は既に同じ無印Jordan類である。
従って全基底群の商ファイバーには他の軌道はない。
\end{proof}
